\section{Building and verifying the formalization}\label{app:verification}

The Lean project is the directory \leanfile{lean/} of this memo. It uses Lean and Mathlib, both at version v4.34.1
(\leanfile{lean-toolchain}, \leanfile{lakefile.toml}, \leanfile{lake-manifest.json}), and requires the parent paper's
library \leanfile{SwapRigidity} through the relative path \leanfile{../../lean}; that library requires the shared
\leanfile{lean-machine/} library of the repository, which supplies Mathlib at the pinned revision. Keep the
repository's directory layout so that both stay in place.

\paragraph{Build.} With \texttt{elan} installed, run in \leanfile{lean/}:
\begin{center}
\texttt{lake exe cache get \&\& lake build}
\end{center}
The first command downloads the compiled Mathlib cache; the second compiles \leanfile{ClassUniformMemo} and, if needed,
the parent paper's library.

\paragraph{Axioms and index.} The file \leanfile{lean/PaperAxioms.lean}, generated from the \LaTeX{} sources of this memo
by \leanfile{tools/lean_index.py}, runs \texttt{\#print axioms} on every declaration cited in this memo
(\cref{app:leanindex}):
\begin{center}
\texttt{lake env lean PaperAxioms.lean}
\end{center}
Every declaration depends at most on \texttt{propext}, \texttt{Classical.choice} and \texttt{Quot.sound}. No Lean file
of this library uses \texttt{sorry}, \texttt{admit} or \texttt{native\_decide}, declares an \texttt{axiom}, bypasses the
kernel via \texttt{unsafe}, \texttt{implemented\_by} or \texttt{extern}, or changes \texttt{maxHeartbeats} or
\texttt{maxRecDepth}. Running \texttt{python3 tools/lean\_index.py} from the memo's root directory regenerates both
\leanfile{paper/sections/leanindex.tex} and \leanfile{lean/PaperAxioms.lean} (\texttt{--check} verifies that every cited
name of this memo's library resolves and that both generated files are current; the names of the parent paper's
library are resolved when \leanfile{PaperAxioms.lean} runs).

\paragraph{Where things are.} The library \leanfile{lean/ClassUniformMemo/} has the following modules, all imported by
the root module \leanfile{lean/ClassUniformMemo.lean}.
\begin{itemize}[leftmargin=1.5em]\raggedright
\item \leanfile{Transfer.lean} (\cref{prop:certificate}): certificates over a commutative ring and their transfer to
fields.
\item \leanfile{RationalData.lean}: the integer data of the two rational codes, generated by
\leanfile{tools/gen_rational_data.py} from \leanfile{tools/data/rational_a.json} and
\leanfile{tools/data/rational_b.json}. It contains definitions only; run the generator with \texttt{--check} to confirm
that it is current.
\item \leanfile{Fields.lean} (\cref{thm:fields}): the three certificates, the finite checks by \texttt{decide}, which
the kernel checks, and the theorem for every field in which $3\ne0$.
\item \leanfile{BlowUp.lean} (\cref{thm:blowup,cor:excess}): blocks, blow-ups, the decoders of a blow-up, and its
dimension.
\item \leanfile{Excess.lean} (\cref{cor:family,prop:merged}): the family on $\Z/6k$, and the bound for the two merged
classes.
\item \leanfile{Advice.lean} (\cref{sec:advice}): linear advice and its exact price, and advice of any kind over a
finite field.
\end{itemize}

\paragraph{Computations.} The only computation is the generation of \leanfile{RationalData.lean}
(\cref{rem:certificates}). The generator needs only the Python standard library. Before it writes the data it re-checks
them, and it also checks the two facts of \cref{rem:certificates} that Lean does not: that the solution space of each
rule has dimension $25$, and that the two codes span a space of dimension $30$. No proof depends on a computation that
Lean does not check.
